
\documentclass{article}
\usepackage{colortbl}
\usepackage{makecell}
\usepackage{multirow}
\usepackage{supertabular}

\begin{document}

\newcounter{utterance}

\centering \large Interaction Transcript for game `cladder', experiment `full\_v1.5\_default', episode 5585 with Qwen3.5{-}9B.
\vspace{24pt}

{ \footnotesize  \setcounter{utterance}{1}
\setlength{\tabcolsep}{0pt}
\begin{supertabular}{c@{$\;$}|p{.15\linewidth}@{}p{.15\linewidth}p{.15\linewidth}p{.15\linewidth}p{.15\linewidth}p{.15\linewidth}}
    \# & \multicolumn{2}{c}{Player} && \multicolumn{2}{c}{Game Master} \\
    \hline

    \theutterance \stepcounter{utterance}  
    & & & \multicolumn{4}{p{0.6\linewidth}}{
        \cellcolor[rgb]{0.9,0.9,0.9}{
            \makecell[{{p{\linewidth}}}]{
                \texttt{\tiny{[P1$\langle$GM]}}
                \texttt{You are an expert in causal inference. The following question is not a typical commonsense query, but rather a meticulously designed question created by a professor specializing in causal inference, intended to assess the students' mastery of the course content.} \\
\texttt{Imagine a self{-}contained, hypothetical world with only the following conditions, and without any unmentioned factors or causal relationships: Gender has a direct effect on treatment and peanut allergy. Treatment has a direct effect on peanut allergy. For patients who are not male and not receiving treatment, the probability of being allergic to peanuts is 84\%. For patients who are not male and receiving treatment, the probability of being allergic to peanuts is 71\%. For patients who are male and not receiving treatment, the probability of being allergic to peanuts is 21\%. For patients who are male and receiving treatment, the probability of being allergic to peanuts is 3\%. The overall probability of male gender is 54\%. For patients receiving treatment, would it be more likely to see being allergic to peanuts if the patient had received no treatment?} \\
\texttt{Start your answer with 'Yes' or 'No', followed by additional reasoning or evidence to support your explanation.} \\
            }
        }
    }
     \\ \\

    \theutterance \stepcounter{utterance}  
    & \multicolumn{4}{p{0.6\linewidth}}{
        \cellcolor[rgb]{0.561,0.737,0.561}{
            \makecell[{{p{\linewidth}}}]{
                \texttt{\tiny{[P1$\rangle$GM]}}
                \texttt{No} \\
\\ 
\texttt{**Reasoning:**} \\
\\ 
\texttt{To answer this question, we must determine the causal effect of treatment on peanut allergy for the specific subgroup of patients who are currently receiving treatment. This requires calculating the **Counterfactual Probability**: \$P(\textbackslash{}text\{Allergy\} | \textbackslash{}text\{Treatment\}=1, \textbackslash{}text\{Gender\}=M)\$ versus \$P(\textbackslash{}text\{Allergy\} | \textbackslash{}text\{Treatment\}=1, \textbackslash{}text\{Gender\}=M)\$, and comparing it to the probability if they had received no treatment, \$P(\textbackslash{}text\{Allergy\} | \textbackslash{}text\{Treatment\}=0, \textbackslash{}text\{Gender\}=M)\$.} \\
\\ 
\texttt{However, the question asks generally: "For patients receiving treatment, would it be more likely to see being allergic to peanuts if the patient had received no treatment?" This implies we need to evaluate the average treatment effect (ATE) or the conditional average treatment effect (CATE) for the treated population.} \\
\\ 
\texttt{Let's break down the provided data into a causal graph structure:} \\
\texttt{*   \$G\$: Gender (Male = 1, Not Male = 0)} \\
\texttt{*   \$T\$: Treatment (Treated = 1, Not Treated = 0)} \\
\texttt{*   \$Y\$: Peanut Allergy (Yes = 1, No = 0)} \\
\\ 
\texttt{**Given Probabilities (Observed Data):**} \\
\texttt{1.  \$P(Y=1 | G=0, T=0) = 0.84\$} \\
            }
        }
    }
    & & \\ \\

    \theutterance \stepcounter{utterance}  
    & & & \multicolumn{2}{p{0.3\linewidth}}{
        \cellcolor[rgb]{0.9,0.9,0.9}{
            \makecell[{{p{\linewidth}}}]{
                \texttt{\tiny{[GM$|$GM]}}
                \texttt{No} \\
            }
        }
    }
    & & \\ \\

    \theutterance \stepcounter{utterance}  
    & & & \multicolumn{2}{p{0.3\linewidth}}{
        \cellcolor[rgb]{0.9,0.9,0.9}{
            \makecell[{{p{\linewidth}}}]{
                \texttt{\tiny{[GM$|$GM]}}
                \texttt{yes} \\
            }
        }
    }
    & & \\ \\

    \theutterance \stepcounter{utterance}  
    & & & \multicolumn{2}{p{0.3\linewidth}}{
        \cellcolor[rgb]{0.9,0.9,0.9}{
            \makecell[{{p{\linewidth}}}]{
                \texttt{\tiny{[GM$|$GM]}}
                \texttt{game\_result = LOSE} \\
            }
        }
    }
    & & \\ \\

\end{supertabular}
}

\end{document}
